\documentclass[11pt]{amsart}
\usepackage{amsmath,amssymb,amsthm}
\usepackage[margin=1.1in]{geometry}
\usepackage{hyperref}
\usepackage[capitalise,nameinlink]{cleveref}
\emergencystretch=2em

\newtheorem{theorem}{Theorem}
\newtheorem{lemma}[theorem]{Lemma}
\newtheorem{corollary}[theorem]{Corollary}
\newtheorem{proposition}[theorem]{Proposition}
\theoremstyle{remark}
\newtheorem{remark}[theorem]{Remark}

\newcommand{\E}{\mathbb{E}}

\title{The stationary cover time of reversible chains on three and four states}
\author{Yitzchak Shmalo}
\date{}

\begin{document}

\begin{abstract}
Aldous and Fill asked whether every reversible Markov chain on $n$ states, started from its stationary law, has
expected cover time at least $(n-1)H_{n-1}$, the value for the simple random walk on the complete graph $K_n$. We
prove this for $n=3$ and $n=4$, for all reversible chains including those with holding, and show that equality holds
only for the simple random walk on $K_n$. The case $n=3$ reduces to Schur's inequality. The case $n=4$ rests on an
explicit polynomial identity of degree $16$ whose numerator becomes coefficientwise nonnegative after an ordered
substitution on each of $30$ cones of edge orderings. Holding is handled for every $n$ by an occupation identity: if a chain
without holding satisfies the bound and $E_iC\ge (n-1)H_{n-1}/2$ for every state $i$, then the bound survives any
added holding. This pointwise condition holds for $n\le 4$.
\end{abstract}

\maketitle

\section{Introduction}

Let $(X_t)_{t\ge0}$ be an irreducible reversible Markov chain on a finite set $V$ of $n$ states, with stationary law
$\pi$. The cover time $C$ is the first time at which every state has been visited; the starting state counts as
visited at time $0$. Open Problem~6.30 in Aldous and Fill~\cite{AF} asks whether
\begin{equation}\label{eq:conj}
\E_\pi C\ \ge\ (n-1)H_{n-1},\qquad H_m=1+\tfrac12+\dots+\tfrac1m,
\end{equation}
for every reversible chain on $n$ states. The right side is the stationary cover time of the simple random walk on
$K_n$: from a uniform start, the walk waits a geometric number of steps with mean $(n-1)/(n-k)$ for the $k$-th new
state, $k=1,\dots,n-1$. For unweighted graphs and the worst starting state, the asymptotic lower bound of order
$n\log n$ is due to Feige~\cite{Feige}; see also~\cite[Chapter~6]{AF}. The exact inequality~\eqref{eq:conj} asks for
more, since it compares every weighted chain, with holding allowed, with one explicit extremal chain. For $n=2$ it is
the statement $\E_\pi C\ge 1$, which is clear.

\begin{theorem}\label{thm:main}
Let $n\in\{3,4\}$. Every irreducible reversible discrete-time Markov chain on $n$ states satisfies
$\E_\pi C\ge (n-1)H_{n-1}$, that is $\E_\pi C\ge 3$ for $n=3$ and $\E_\pi C\ge 11/2$ for $n=4$. Equality holds if and
only if the chain is the simple random walk on $K_n$: every off-diagonal transition probability is $1/(n-1)$ and
there is no holding.
\end{theorem}

The proof has three parts. \Cref{sec:formula} writes $\E_\pi C$ as an alternating sum of quadratic forms in inverses
of principal submatrices of the weighted Laplacian; this holds for every $n$. \Cref{sec:holding} shows that holding
can be added afterwards, under a pointwise condition that is automatic for $n\le4$
(\cref{prop:holding}). Holding can lower $\E_\pi C$ for a fixed jump chain (a three-state example is given in
\cref{rem:holding-lowers}), so the reduction is not a monotonicity statement. \Cref{sec:three,sec:four} prove the
bound without holding for $n=3$ and $n=4$.

\section{The stationary cover time as an alternating sum}\label{sec:formula}

A reversible chain on $V$ is described by symmetric conductances $w_{xy}=w_{yx}\ge0$ for $x\ne y$ and holding weights
$h_x\ge0$: with $q_x=\sum_{y\ne x}w_{xy}$ and $m_x=q_x+h_x$, the chain moves from $x$ to $y\ne x$ with probability
$w_{xy}/m_x$ and stays at $x$ with probability $h_x/m_x$. Its stationary law is $\pi_x=m_x/\sum_ym_y$. Every
irreducible reversible chain arises this way, with $w_{xy}=\pi_xP_{xy}$ and $h_x=\pi_xP_{xx}$, and irreducibility means
that the graph of positive conductances is connected. In this section and in \cref{sec:three,sec:four} there is no
holding, so $m=q$, and we write $W=\sum_xq_x$.

Let $L$ be the weighted Laplacian, $L_{xx}=q_x$ and $L_{xy}=-w_{xy}$ for $x\ne y$. For a proper nonempty subset
$S\subset V$, let $L_S$ be the principal submatrix on $S$ and $q_S=(q_x)_{x\in S}$. If a vector $v$ supported on $S$
has $v^{\top}Lv=\tfrac12\sum_{x,y}w_{xy}(v_x-v_y)^2=0$, it is constant on the connected support and vanishes off $S$,
so $v=0$. Hence $L_S$ is positive definite.

\begin{lemma}\label{lem:hitting}
For a nonempty set $J\subset V$ with complement $S\ne\emptyset$, let $T_J$ be the hitting time of $J$. Then
$\E_\pi T_J=q_S^{\top}L_S^{-1}q_S/W$.
\end{lemma}

\begin{proof}
The function $g_x=\E_xT_J$ vanishes on $J$ and satisfies $g_x=1+\sum_yP_{xy}g_y$ for $x\in S$. Multiplying by $q_x$
gives $q_xg_x-\sum_{y\in S}w_{xy}g_y=q_x$, that is $L_Sg_S=q_S$. Averaging against $\pi_x=q_x/W$ gives the claim.
\end{proof}

Let $T_x$ be the hitting time of the single state $x$, with $T_x=0$ if $X_0=x$. Then $C=\max_xT_x$. For real numbers
$t_x$, $\max_xt_x=\sum_{\emptyset\ne J\subseteq V}(-1)^{|J|+1}\min_{x\in J}t_x$, and $\min_{x\in J}T_x=T_J$. Taking
expectations, using \cref{lem:hitting}, and writing $S=V\setminus J$ (the term $J=V$ vanishes),
\begin{equation}\label{eq:maxmin}
W\,\E_\pi C=\sum_{\emptyset\ne S\subsetneq V}(-1)^{n-|S|+1}\,q_S^{\top}L_S^{-1}q_S .
\end{equation}

\section{Adding holding}\label{sec:holding}

\begin{proposition}\label{prop:holding}
Let $Q$ be an irreducible reversible chain on $n$ states without holding, with conductances $w$, and let $c_0>0$.
Suppose that $\E_\pi C\ge c_0$ for $Q$ and that $\E_xC\ge c_0/2$ for every state $x$. Then every chain obtained from
$Q$ by adding holding weights $h\ge0$ satisfies $\E_{\pi_h}C_h\ge c_0$. More precisely, with $W_h=W+\sum_xh_x$,
\begin{equation}\label{eq:holding}
W_h\bigl(\E_{\pi_h}C_h-c_0\bigr)=W\bigl(\E_\pi C-c_0\bigr)+\sum_xh_x\bigl(2\,\E_xC-c_0\bigr)+h^{\top}Kh
\end{equation}
for a symmetric matrix $K$ with nonnegative entries.
\end{proposition}

\begin{proof}
For the chain $Q$ let $N_{xy}=\E_x\sum_{t=0}^{C-1}\mathbf 1\{X_t=y\}$ and $K_{xy}=N_{xy}/q_y$. For each $t$, the event
$\{t<C\}$ depends only on the path $X_0,\dots,X_t$, through the condition that it misses some state, and this condition
is unchanged when the path is reversed. Reversibility gives $q_x\,\mathbb P(\text{path})=q_y\,\mathbb P(\text{reversed
path})$ for a path from $x$ to $y$. Summing over paths of length $t$ that miss a state, and then over $t$, gives
$q_xN_{xy}=q_yN_{yx}$, so $K$ is symmetric. Since $\sum_yN_{xy}=\E_xC$, we have $(Kq)_x=\E_xC$.

Add holding weights $h$ and let $m=q+h$. The jump chain of the new chain is $Q$, and each visit to $y$ before the last
new state is reached lasts a geometric number of steps with mean $m_y/q_y$. The cover time ends on arrival at the last
new state, so no waiting is counted there. Conditioning on the jump chain gives $\E_xC_h=\sum_yK_{xy}m_y$, and with
$\pi_h(x)=m_x/W_h$,
\[
W_h\,\E_{\pi_h}C_h=m^{\top}Km=q^{\top}Kq+2h^{\top}Kq+h^{\top}Kh=W\,\E_\pi C+2\sum_xh_x\,\E_xC+h^{\top}Kh .
\]
Subtracting $c_0W_h=c_0W+c_0\sum_xh_x$ gives~\eqref{eq:holding}. Every term on the right of~\eqref{eq:holding} is
nonnegative under the hypotheses.
\end{proof}

Any chain needs at least $n-1$ steps to cover $n$ states, so $\E_xC\ge n-1$. For $c_0=(n-1)H_{n-1}$ the condition
$n-1\ge c_0/2$ reads $H_{n-1}\le2$, which holds for $n\le4$ (with $H_3=11/6$) and fails for $n\ge5$. For $n=3$ and
$n=4$ the middle term of~\eqref{eq:holding} is at least $\sum_xh_x$ and $\tfrac12\sum_xh_x$ respectively, so any
positive holding gives a strict inequality. Hence \cref{thm:main} follows from its case without holding, which is
proved in the next two sections.

\begin{remark}\label{rem:holding-lowers}
Adding holding to a fixed jump chain can decrease $\E_\pi C$, because holding also changes the stationary law, and
hence the starting distribution. Take three states with $w_{12}=1$ and $w_{13}=w_{23}=1/100$, so that state $3$ is
hard to reach. Without holding $\E_\pi C\approx100.04$; with holding weight $h_3=1$ at state $3$ the chain starts
there more often, and $\E_\pi C\approx85.07$. \Cref{prop:holding} therefore cannot be replaced by the statement that
holding never lowers the stationary cover time.
\end{remark}

\section{Three states}\label{sec:three}

Write $a=w_{12}$, $b=w_{13}$, $c=w_{23}$, with at most one of them zero. By~\eqref{eq:maxmin}, or by a direct
first-step computation,
\begin{equation}\label{eq:three}
\E_\pi C-3=\frac{Q(a,b,c)}{2(a+b+c)(ab+bc+ca)},\qquad
Q=4(a^3+b^3+c^3)-\sum_{\mathrm{sym}}a^2b-6abc,
\end{equation}
where $\sum_{\mathrm{sym}}a^2b$ is the sum of the six monomials $a^2b,a^2c,b^2a,b^2c,c^2a,c^2b$. Now
\[
Q=\Bigl(a^3+b^3+c^3+3abc-\sum_{\mathrm{sym}}a^2b\Bigr)+3\bigl(a^3+b^3+c^3-3abc\bigr).
\]
The first bracket is nonnegative by Schur's inequality and the second by the inequality of arithmetic and geometric
means. The second vanishes only when $a=b=c$, and connectivity then forces $a=b=c>0$. Hence $\E_\pi C\ge3$ without
holding, with equality only for the simple random walk on $K_3$. Together with \cref{prop:holding} this proves
\cref{thm:main} for $n=3$.

\section{Four states}\label{sec:four}

Label the edges of $K_4$ as $(x_0,\dots,x_5)=(w_{01},w_{02},w_{03},w_{12},w_{13},w_{23})$. For the sets $S$ of size
one, $q_S^{\top}L_S^{-1}q_S=q_x$, and these terms sum to $W$. For $S=\{i,j\}$,
\[
q_S^{\top}L_S^{-1}q_S=\frac{B_{ij}}{\Delta_{ij}},\qquad \Delta_{ij}=q_iq_j-w_{ij}^2,\qquad
B_{ij}=q_iq_j(q_i+q_j+2w_{ij}).
\]
For $|S|=3$, the determinant of $L_S$ is the spanning-tree polynomial $\tau$ of the weighted $K_4$, by the matrix-tree
theorem. Put $A=\sum_{|S|=3}q_S^{\top}\operatorname{adj}(L_S)\,q_S$ and $R=\prod_{i<j}\Delta_{ij}$. Then
\eqref{eq:maxmin} becomes $\E_\pi C=1+\bigl(A/\tau-\sum_{i<j}B_{ij}/\Delta_{ij}\bigr)/W$, so
\begin{equation}\label{eq:F}
\E_\pi C-\frac{11}{2}=\frac{F}{2W\tau R},\qquad
F=2AR-2\tau\sum_{i<j}B_{ij}\prod_{\{k,l\}\ne\{i,j\}}\Delta_{kl}-9W\tau R .
\end{equation}
The denominator is positive: $\tau$ and $\Delta_{ij}$ are determinants of the positive definite matrices $L_S$. The
numerator $F$ is a homogeneous polynomial of degree $16$ in $x_0,\dots,x_5$ with integer coefficients. It has
$13{,}368$ terms, of which $2{,}421$ have negative coefficients, so its nonnegativity is not visible term by term.

\begin{proposition}\label{prop:cones}
$F(x)\ge0$ for every $x\in[0,\infty)^6$ whose support is a connected spanning subgraph of $K_4$, with equality only
when all six $x_k$ are equal.
\end{proposition}

\begin{proof}
Every $x\in[0,\infty)^6$ has some ordering $x_{o_0}\le x_{o_1}\le\dots\le x_{o_5}$, where $o$ is a permutation of the
six edges. Writing $x_{o_k}=z_0+z_1+\dots+z_k$ with $z_k\ge0$ turns $F$ into a polynomial $G_o(z)$. Relabelling the
four states permutes the edges and leaves $F$ unchanged, since it leaves $\E_\pi C$, $W$, $\tau$ and $R$ unchanged. The
$720$ orderings fall into $30$ orbits of size $24$ under the $24$ relabellings, because no relabelling other than the
identity fixes all six edges. It is therefore enough to treat one representative $o$ of each orbit.

For each of the $30$ representatives, the polynomial $G_o$ has nonnegative integer coefficients; there are
$577{,}620$ nonzero coefficients in all, recorded in the accompanying files. This proves $F\ge0$. For equality, note that
in every $G_o$ the monomials $z_0^{14}z_j^2$, $1\le j\le5$, have positive coefficients. If all edges are positive then
$z_0>0$, so $F=0$ forces $z_1=\dots=z_5=0$, which means all six weights are equal. If some edge vanishes, let $j\ge1$ be
the smallest index with $x_{o_j}>0$; then $z_j>0$, and the positive edges are exactly $o_j,\dots,o_5$. Whenever these
edges form a connected spanning subgraph, the coefficient of $z_j^{16}$ in $G_o$ is positive, so $F>0$.
\end{proof}

By~\eqref{eq:F} and \cref{prop:cones}, $\E_\pi C\ge11/2$ for every irreducible reversible chain on four states without
holding, with equality only when all conductances are equal, that is for the simple random walk on $K_4$. With
\cref{prop:holding} this proves \cref{thm:main} for $n=4$.

\begin{remark}
\Cref{prop:cones} is a finite computation. The coefficients of the $30$ polynomials $G_o$, the polynomial $F$, and
programs that rebuild $F$ from~\eqref{eq:maxmin} and check every coefficient in exact integer arithmetic are in the
repository listed below. Two further checks are independent of the program that produced the certificate: a second
reconstruction of $F$ from integer dictionaries, Leibniz determinants and an enumeration of spanning trees, and an
evaluation of each $G_o$ at random integer points against $2W\tau R(\E_\pi C-11/2)$ computed by an exact solve of the
first-step equations on pairs (current state, visited set). Exact enumeration over all $624$ connected weightings in
$\{0,1,2\}^6$, each with three holding vectors, finds minimum $11/2$, attained only at the two constant weightings
without holding. The
case $n=3$ with holding, and \cref{prop:holding} in the form of a statement about symmetric matrices, have been checked in
Lean~4 with Mathlib.
\end{remark}

\section{Remarks on larger $n$}

For $n\ge5$ the trivial bound $\E_xC\ge n-1$ falls short of $(n-1)H_{n-1}/2$, so the hypothesis of
\cref{prop:holding} would have to be proved separately, and the reduction from chains with holding is open. The approach
of \cref{sec:four} also grows quickly: for $n=5$ there are ten edges and $10!/120=30{,}240$ orbits of edge orderings.
Numerical minimisation over reversible chains on five and six states, with holding, has found no value below
$(n-1)H_{n-1}$.

\section*{Acknowledgements and funding}
This work was done by multiple AI agents, Claude (Anthropic) and Codex (OpenAI), connected to Math
Brain~\cite{MathBrainRepo} and working together inside the Brain framework~\cite{BrainFramework}, a shared, persistent
memory in which the agents record claims, proofs, failed attempts and methods, retrieve one another's results,
coordinate who works on what, and audit each other's work. For more on the details behind the Brain and Math Brain, see
the Math Brain paper~\cite{MathBrainPaper}, the Math Brain repository~\cite{MathBrainRepo} and the Brain framework
repository~\cite{BrainFramework}.

The research presented in this paper was supported by the European Research Council (ERC) under the European Union's
Horizon Europe research and innovation programme (grant agreement No.~101041711), by the Simons Foundation as part of
the Collaboration on the Mathematical and Scientific Foundations of Deep Learning, by Heights Labs, by Convex Nexus
Capital, by the Israel Science Foundation (grant number 2258/19), and by the Israel Science Foundation (ISF Grant
4101/25). Computational resources were provided in part by Google Cloud through its education credits program (credit
530550103).

\section*{Declaration of competing interest}
The author holds equity in Heights Labs and in Convex Nexus Capital and has a management role at Convex Nexus Capital.
These entities are included in the funding declaration above.

\section*{Declaration of generative AI and AI-assisted technologies}
The paper comes out of Math Brain, an ongoing project on the use of AI systems for research mathematics, and its
results were developed with the Brain~\cite{BrainFramework} and Math Brain~\cite{MathBrainRepo}. In October and
November 2026 Claude agents (Anthropic) and Codex agents running OpenAI's GPT-6.1, connected to the Brain and working
on the same problem, studied the stationary cover time of small reversible chains. A Claude agent proved the case
$n=3$; a Codex agent found the formula~\eqref{eq:F}, the cone certificate for $n=4$ and the occupation identity of
\cref{prop:holding}. Claude agents checked the certificate with independent programs, formalised the case $n=3$ and the
holding lemma in Lean, and Claude Opus 5.5 (Anthropic) prepared the manuscript and audited it. These checks are not peer
review. The author directed the work and is responsible for the content of the paper and for decisions concerning its
dissemination.

\section*{Code and data availability}
Code, certificates and Lean files are at

\url{https://github.com/yspennstate/stationary-cover-time-four-states}.

The paper is part of the Math Brain project, whose public repository is
\url{https://github.com/yspennstate/math-brain-public}~\cite{MathBrainRepo}.

\begin{thebibliography}{99}
\bibitem{AF} D.~Aldous, J.~A.~Fill, \emph{Reversible Markov Chains and Random Walks on Graphs}, unfinished monograph,
recompiled 2014, \url{https://www.stat.berkeley.edu/~aldous/RWG/book.html}.
\bibitem{Feige} U.~Feige, A tight lower bound on the cover time for random walks on graphs, Random Structures
Algorithms 6 (1995), 433--438.
\bibitem{MathBrainRepo} Y.~Shmalo, \emph{Math Brain: public snapshot, code and tests}, 2026,
\url{https://github.com/yspennstate/math-brain-public}.
\bibitem{MathBrainPaper} Y.~Shmalo, \emph{Math Brain: shared, auditable memory for multi-agent mathematical research},
October 2026, \url{https://github.com/yspennstate/math-brain-public/tree/main/papers/math-brain-2026-10-24}.
\bibitem{BrainFramework} Y.~Shmalo, \emph{Brain Framework}, public software repository, 2026,
\url{https://github.com/yspennstate/brain-framework-public}.
\end{thebibliography}

\end{document}
